\documentclass{article}
% Each code line below has four leading spaces removed by gobble.
\usepackage[gobble=4,tabsize=4]{paper-utils}

\title{A Brief Guide to the \texttt{verbatim} package}
\date{\today}
\author{\href{https://pages.di.unipi.it/massa}{Jacopo Massa}}

\begin{document}

\maketitle

\section*{Options}
Be sure to set the options \verb|tabsize| and \verb|gobble| in the preamble, as shown below.

\begin{Verbatim}
	\usepackage[tabsize=4,gobble=1]{paper-utils}
\end{Verbatim}

\noindent
\verb|tabsize| (default: 4) sets the number of spaces per tab character.\\
\verb|gobble| (default: 4) sets the number of characters to remove from the beginning of each line of code. This is useful for removing leading whitespace in indented code blocks.\\

\noindent In the \texttt{paper-utils} package, there are two environments for including code: \verb|code| (for unnumbered code blocks) and \verb|codenum| (for numbered ones).

\section*{\texttt{code} environment}
\begin{code}
    application(AppId, [Functions], [Services]).
\end{code}

\section*{\texttt{codenum} environment}
\begin{codenum}[firstnumber=1]
    fogBrain(App, NewPlacement) :-
        deployment(App, Placement, AllocHW, AllocBW),
        reasoningStep(App, Placement, AllocHW, AllocBW, NewPlacement).

    fogBrain(App, Placement) :-
        placement(App, Placement).
\end{codenum}

\section*{Comments}
\begin{codenum}[firstnumber=1]
    application(app, [], []). \verbcomment{% Example application}
\end{codenum}

\section*{``\texttt{.localSettings.yaml}'' file}

If you use \cd{latexIndent} on your local machine, you can copy the \cd{.localSettings.yaml} file from the \cd{packages} folder to the directory where your \cd{.tex} files are located. This will deactivate the auto-indentation of the \cd{code} and \cd{codenum} environments.

\end{document}
